\label{sec:related}

The closest literature separates into finite-element automation,
differentiable mechanics, constitutive differentiation, and sparse derivative
recovery. We organize it by the location of differentiation because software
feature inventories do not determine whether two routes represent the same
discrete functional.

The FEniCS project established symbolic weak-form specification and generated
finite-element kernels \citep{logg2012fenicsbook}. DOLFINx extends this model
with data-oriented parallel interfaces \citep{baratta2025dolfinx}. Automated
adjoint systems operate at a higher program level: \code{dolfin-adjoint} and
\code{pyadjoint} derive sensitivities of finite-element programs
\citep{farrell2013dolfinadjoint,mitusch2019pyadjoint}, while \code{cashocs}
provides adjoint-based optimal-control and shape-optimization workflows
\citep{blauth2023cashocsv2}. These systems motivate automation, but the present
work concerns primal residual and tangent construction inside a Newton solve.

JAX provides composable program transformations for differentiation,
vectorization, and compilation \citep{jax2018,baydin2018autodiff}. JAX-FEM
applies them to forward and inverse finite-element mechanics
\citep{xue2023jaxfem}, and second-order implicit differentiation extends this
setting to Hessian-vector products for differentiable physics
\citep{xue2026implicit}. AutoPDEx and JAX-CPFEM illustrate related
energy-based and constitutive applications
\citep{bode2025autopdex,hu2025jaxcpfem}. These results establish the viability
of differentiable finite elements; they do not compare the three derivative
locations studied here under one distributed algebraic contract.

Bridge architectures are particularly close. Firedrake--JAX differentiates
partial differential equation solves through tangent and adjoint equations
\citep{yashchuk2023bringing}.
FEniCSx external operators admit general quadrature-point programs, including
a JAX-differentiated Mohr--Coulomb model with Taylor tests
\citep{latyshev2025externaloperators}. JetSCI combines JAX-local
discretizations with PETSc-owned distributed solves
\citep{cattaneo2026jetsci}. Thus neither constitutive AD nor the JAX--PETSc
architecture is new here. Our distinction is the fixed-functional equivalence
result and an evidence hierarchy for comparing element, constitutive, and
colored routes.

Recent systems narrow this distinction further. Locality-aware AD performs
per-element forward differentiation and sparse Hessian assembly on graphics
processing units
\citep{mahmoud2025locality}. The \code{tatva} framework applies AD to a global
energy and uses matrix-free products or coloring to materialize sparse
tangents \citep{pundir2026tatva}. Hill and Dalle combine automatic local or
global sparsity detection, matrix coloring, and sparse differentiation
\citep{hill2025sparser}. These papers establish both local and global AD
strategies. The present contribution is therefore limited to explicit
branchwise equivalence assumptions, distributed finite-element ownership
conditions, and a fixed-state verification hierarchy on processor-based MPI
systems.

Sparse recovery itself is classical. Curtis, Powell, and Reid introduced
directional probing for sparse Jacobians \citep{curtis1974sparsejacobian};
Coleman and Mor{\'e} related sparse Jacobian and Hessian estimation to graph
coloring \citep{coleman1983sparsejacobian,coleman1984sparsehessian}. PETSc
supplies the distributed vectors, matrices, Krylov methods, and preconditioners
used in our implementation \citep{balay1997petsc,petsc2026manual}. We do not
claim these algorithms as new. We use them to expose the assumptions needed to
turn local Hessian actions into a canonical distributed matrix and to separate
finite correctness evidence from performance claims.

The branch-structured benchmark is motivated by Mohr--Coulomb slope stability,
but its evidentiary role is intentionally narrower. Incremental plasticity,
return mapping, and consistent linearization provide the continuum reference
frame \citep{simo1985consistent,simo1998compinel,sysala2017returnmapping}.
Variational strength reduction and recent three-dimensional continuation
methods provide application context
\citep{sysala2021optimization,sysala2025advancedcontinuation}. Our implemented
object is a synthetic endpoint potential with selected branches. It tests
derivative placement; it does not validate a path-consistent constitutive
history or a factor of safety.
